\chapter{Compiler architecture overview (prototype + Babel transform)}
\label{ch:compiler-architecture}

This chapter explains how Relax's HRBR compiler takes TSX/JSX input and produces output that mounts \emph{compiled blocks} (the fast path) or routes to the \emph{fallback runtime} when the template's structure is dynamic.

\subsection*{What you\'ll build (mentally)}

This chapter is about building a reliable mental model, not shipping a production compiler.
By the end, you should be able to:

\begin{itemize}
\item explain the compiler/runtime boundary in one sentence
\item read transformed output and understand where \texttt{templateHTML}, \texttt{slots}, and \texttt{read()} closures come from
\item predict when compilation must route to \texttt{mountFallback()} because the DOM structure isn\'t stable
\item debug slot path bugs by reasoning about \texttt{childNodes} indexing (elements \emph{and} text nodes)
\end{itemize}

\subsection*{Prerequisites}

You\'ll get the most value if you already understand:

\begin{itemize}
\item HRBR blocks: \texttt{defineBlock()} + \texttt{mountCompiledBlock()} (Chapter~\ref{ch:hrbr-blocks-template-slots})
\item fallback reconciliation: \texttt{mountFallback()} + structural patching (Chapter~\ref{ch:hrbr-fallback-reconciler})
\item VDOM \(\leftrightarrow\) HRBR interop boundary: \texttt{hBlock(mount)} (Chapter~\ref{ch:vdom-hrbr-interop})
\end{itemize}

\subsection*{Why this matters}

Compilers are optional in UI libraries until they aren\'t.
The moment you want both:

\begin{itemize}
\item a stable template that never re-creates DOM, and
\item ergonomic authoring in JSX/TSX,
\end{itemize}

you end up building a compiler or transform.

This chapter is deliberately practical: we\'ll treat tests as executable contracts, and we\'ll keep the compiler/runtime boundary small enough that you can maintain it.

The compiler currently exists in two forms:

\begin{itemize}
\item A TypeScript-AST prototype compiler in \texttt{compiler/index.ts} (used as an executable spec for template extraction).
\item A Babel plugin scaffold in \texttt{compiler/jsx/transform.ts} (used by build tooling to transform real TSX/JSX files).
\end{itemize}

The purpose of having both is pedagogical and practical:

\begin{itemize}
\item The TS compiler is a small, deterministic ``reference'' that focuses on \emph{template + slot-path extraction}.
\item The Babel plugin is the real integration point for user codebases.
\end{itemize}

\section{The output contract (what the runtime needs)}

Before we talk about parsing and visitors, it's worth making the compiler/runtime boundary explicit.
The compiler's job is to emit \emph{data and closures} that satisfy a small runtime API.
Everything in the implementation is organized around making that boundary stable.

\subsection*{Diagram: the compiler-to-runtime boundary (block vs fallback)}
\begin{center}
\begin{tikzpicture}[
	node distance=10mm and 14mm,
	box/.style={draw, rounded corners, align=left, inner sep=6pt, fill=RelaxLight},
	arrow/.style={-Latex, thick}
]
\node[box] (jsx) {Input TSX/JSX\\\texttt{<div>...}};
\node[box, right=of jsx] (compiler) {Compiler\\(TS prototype or Babel plugin)};
\node[box, right=of compiler] (decision) {Route decision\\stable structure?};

\node[box, above right=of decision] (block) {Block mode\\\texttt{defineBlock(\{templateHTML, slots\})}\\\texttt{compiledSlots: \{key, read\}[]}\;};
\node[box, below right=of decision] (fallback) {Fallback mode\\emit mount factory\\\texttt{mountFallback(host, renderFn)}};

\node[box, right=of block] (runtimeB) {HRBR runtime\\\texttt{mountCompiledBlock(def, host, compiledSlots)}};
\node[box, right=of fallback] (runtimeF) {Fallback reconciler\\(runtime)};

\draw[arrow] (jsx) -- (compiler);
\draw[arrow] (compiler) -- (decision);
\draw[arrow] (decision) -- node[midway, above] {yes} (block);
\draw[arrow] (decision) -- node[midway, below] {no} (fallback);
\draw[arrow] (block) -- (runtimeB);
\draw[arrow] (fallback) -- (runtimeF);
\end{tikzpicture}
\end{center}

This is the architectural backbone: the compiler emits either a \emph{BlockDef + read closures} (fast) or a \emph{fallback mount factory} (safe).

\subsection{Tiny contract (2 bullets)}

Given a JSX expression, the compiler must be able to produce one of:

\begin{itemize}
\item \textbf{Block mode (fast path):} \texttt{defineBlock({ templateHTML, slots })} plus a mount factory that calls \texttt{mountCompiledBlock(def, host, compiledSlots)}.
\item \textbf{Fallback mode (safe path):} a mount factory that calls \texttt{mountFallback(...)} and delegates structural reconciliation to the runtime.
\end{itemize}

The runtime intentionally does not know how compilation works; it only requires these shapes.

The HRBR runtime expects two main pieces of data.

\subsection{1) A \texttt{BlockDef}: stable template + slot metadata}

Defined in \texttt{runtime/block.ts}, a \texttt{BlockDef} contains:

\begin{itemize}
\item \texttt{templateHTML: string}, the stable HTML string that will be parsed once.
\item \texttt{slots: Record<string, BlockSlot>} , the mapping from a slot key to its write kind and a DOM path.
\end{itemize}

A \texttt{BlockSlot} describes \emph{what} to write (text/attr/class/style/prop/event, etc.) and \emph{where} to write it (a childNodes index path).

\subsection{2) A \texttt{CompiledSlot[]} list: how to read live values}

At runtime, \texttt{mountCompiledBlock(def, host, compiledSlots)} needs an array of objects shaped like:

\begin{quote}
\texttt{\{ key: string, read: () => unknown \}}
\end{quote}

The compiler can't always inline the read closures in a pure data file.
In particular, the TypeScript prototype compiler deliberately returns no reads; it records expression source code so a caller with lexical scope can build closures.

In contrast, the Babel transform \emph{can} emit \texttt{read: () => expr} arrow functions directly into the output program, precisely because it is already transforming a file with the correct lexical scope.

\subsection{A 3-part boundary you should memorize}

At the boundary, the compiler emits three things:

\begin{enumerate}
\item a \texttt{BlockDef} (pure data: stable HTML + slot metadata)
\item a \texttt{compiledSlots[]} list (closures: how to read live values)
\item a mount factory (a closure boundary: where lexical scope lives)
\end{enumerate}

That last bullet is easy to miss.
It\'s the reason the Babel transform can emit \texttt{read: () => expr} directly, while the TypeScript prototype compiler intentionally stops at ``expression source code''.

For a background overview of why JSX is typically compiled, see:\
\href{https://react.dev/learn/writing-markup-with-jsx}{https://react.dev/learn/writing-markup-with-jsx}.

\section{Routing decision: block vs fallback}

The compiler's core architectural decision is: \emph{can this TSX be represented as a stable HTML template?}
If not, we must use the fallback reconciler.

\subsection{Dynamic structure}

The term \textbf{dynamic structure} means ``the set/shape/order of DOM nodes can change as state changes.''
Examples:

\begin{itemize}
\item conditionals: \texttt{cond ? <div/> : null}
\item loops/lists: \texttt{items.map(...)}
\item fragments that change child count
\item dynamic tag names or component tags
\item spreading attributes (often implies unknown keys)
\item event handlers (in the current compiler phase) if they can't be represented as static bindings
\end{itemize}

Compiled blocks are designed for \emph{stable structure} with dynamic \emph{values}.
When structure is dynamic, we let the fallback runtime reconcile children.

\subsection*{Diagram: what counts as “dynamic structure” (conservative)}
\begin{center}
\begin{tikzpicture}[
	node distance=8mm and 12mm,
	box/.style={draw, rounded corners, align=left, inner sep=6pt, fill=RelaxLight},
	arrow/.style={-Latex, thick},
	decision/.style={draw, diamond, aspect=2, inner sep=2.5pt, align=center, fill=RelaxLight}
]
\node[box] (expr) {JSX expression\\\texttt{<div>...</div>}};
\node[decision, right=of expr] (dyn) {Any child\\expression can\\change node shape?};
\node[box, above right=of dyn] (fallback) {Route to fallback\\(reconcile structure)};
\node[box, below right=of dyn] (block) {Allow block mode\\(write values only)};

\draw[arrow] (expr) -- (dyn);
\draw[arrow] (dyn) -- node[midway, above] {yes} (fallback);
\draw[arrow] (dyn) -- node[midway, below] {no} (block);
\end{tikzpicture}
\end{center}

Relax’s current compilers deliberately err on the safe side: anything that could introduce or remove DOM nodes is treated as “dynamic structure”.

\subsection{A useful rule of thumb}

If an expression could evaluate to:

\begin{itemize}
\item \texttt{null}
\item \texttt{false}
\item an array of nodes
\item a \texttt{<div>} sometimes and \texttt{<span>} other times
\end{itemize}

then it’s structural and must go to fallback.

If it always evaluates to a \emph{value} that can be written into an existing node (text/attr/prop/class/style), then it can stay in block mode.

This is conservative, but it keeps the compiler honest.

\subsection{How ``block vs fallback'' relates to Chapter~\ref{ch:vdom-hrbr-interop}}

The compiler output doesn't mount directly into the VDOM. Instead, it produces a \emph{mount factory} that matches the HRBR vnode contract described in Chapter~\ref{ch:vdom-hrbr-interop}.
That is why the VDOM type \texttt{DOM\_TYPES.HRBR} stores a \texttt{mount(host)} function: it is the common denominator between compiled blocks and fallback regions.

\section{TypeScript prototype compiler (Phase 5 reference)}

The TypeScript implementation lives in \texttt{compiler/index.ts}.
It provides two public APIs:

\begin{itemize}
\item \texttt{compileTSXToBlock(source, opts?) -> CompileResult}
\item \texttt{compileTSXOrFallback(source, opts?) -> CompileOrFallbackResult}
\end{itemize}

\subsection{\texttt{compileTSXToBlock():} parse -> extract -> render template}

\textbf{Inputs}

\begin{itemize}
\item \texttt{source: string}: a single-root TSX expression (or a statement containing one).
\item \texttt{opts.fileName?}: helps TypeScript provide useful source locations.
\item \texttt{opts.strictTemplate?} (default true): reject patterns that can't be expressed as a static HTML template.
\end{itemize}

\textbf{Main pipeline}

\begin{enumerate}
\item Parse the input with TypeScript: \texttt{ts.createSourceFile(..., ScriptKind.TSX)}.
\item Find the first expression via \texttt{findFirstExpression()}.
\item Ensure the root is \texttt{JsxElement} or \texttt{JsxSelfClosingElement}.
\item Build a \texttt{TemplateBuild} using \texttt{buildTemplate()}.
\item Return \texttt{\{ block: \{templateHTML, slots\}, slots: [], meta \}}.
\end{enumerate}

\subsection{Determinism: why the TS compiler is a good executable spec}

The implementation in \texttt{compiler/index.ts} is deliberately small and deterministic:

\begin{itemize}
\item It parses with TypeScript's parser (so syntax + TSX rules match what users write).
\item It walks the syntax tree into a simple string renderer that produces \texttt{templateHTML}.
\item It computes slot paths using an explicit \texttt{childIndex} counter.
\end{itemize}

This makes it ideal for teaching and for pinning down invariants (especially path semantics) before committing to a more complex Babel visitor implementation.

\subsection{Paths: the compiler and runtime must agree}

The most important shared invariant between compiler and runtime is how a ``path'' is interpreted.

In this prototype compiler, each element maintains a \texttt{childIndex} counter, and each child node increments the counter as it is rendered.
When the compiler emits a slot, it records the current path (a list of child indices) that identifies a target node under the template root.

\paragraph{Why this works}

The block runtime later parses \texttt{templateHTML} into a DOM tree and walks \texttt{childNodes[index]} repeatedly to locate the slot target.

\subsection{The one gotcha that causes most bugs: \texttt{childNodes} includes text nodes}

If you take nothing else from this chapter, take this:

\begin{quote}
	extbf{Slot paths are based on DOM \texttt{childNodes}, not \texttt{children}.}
\end{quote}

That means whitespace and literal text in your template count as nodes.
So when you look at markup like:

\begin{quote}
	exttt{<p>Hello \{name()\}</p>}
\end{quote}

the compiled \texttt{templateHTML} might contain literal spaces or newlines, and those become \texttt{Text} nodes when parsed.

This is also why the tests talk about paths like \texttt{[3, 1]} (not \texttt{[0, 0]}).
In the Babel snapshot test (\texttt{compiler/\_\_tests\_\_/babel-transform.test.ts}), the template contains indentation and newlines.
When parsed, they become text nodes that affect indexing.

\subsection{Walkthrough: compute a slot path by hand}

Let\'s do a small thought experiment that matches the compiler\'s rules.

Input TSX:

\begin{lstlisting}[language=JavaScript]
const App = () => (
	<section id="root" className={cls()}>
		<h1>Title</h1>
		<p>Hello {name()}</p>
	</section>
)
\end{lstlisting}

The Babel transform snapshot illustrates this exact structure as a static template plus two slots:

\begin{itemize}
\item \texttt{s0}: a \texttt{class} slot on the root element, with \texttt{path: []}
\item \texttt{s1}: a \texttt{text} slot inside the \texttt{<p>}, with \texttt{path: [3, 1]}
\end{itemize}

Why \texttt{[3, 1]}?

\begin{itemize}
\item \texttt{[]} means \"start at the template root element\" (the \texttt{<section>}).
\item \texttt{3} means \"go to the 4th \texttt{childNodes} entry under \texttt{<section>}\".
	In formatted templates, the \texttt{childNodes} often look like: newline text, \texttt{<h1>}, newline text, \texttt{<p>}, newline text.
\item \texttt{1} then means \"go to the 2nd \texttt{childNodes} entry under \texttt{<p>}\".
		Under \texttt{<p>}, you often have: \texttt{Text('Hello ')} then \texttt{Text('\_\_hrbr\_\_')} (or an empty placeholder) depending on how the compiler encodes it.
\end{itemize}

This is ugly but reliable.
And reliability is the whole point: the compiler and runtime just need to agree.

For background on DOM node lists, see MDN:
\href{https://developer.mozilla.org/en-US/docs/Web/API/Node/childNodes}{https://developer.mozilla.org/en-US/docs/Web/API/Node/childNodes}.

\subsection{What the TS prototype actually returns (and why the reads are empty)}

In \texttt{compiler/index.ts}, \texttt{compileTSXToBlock()} returns:

\begin{itemize}
\item \texttt{block: \{ templateHTML, slots \}}
\item \texttt{slots: []} (yes, empty)
\item \texttt{meta} about dynamic structure, lanes, etc.
\end{itemize}

That\'s not a bug.
It\'s an explicit constraint: a string-only compiler doesn\'t have access to the lexical scope where \texttt{count()}, \texttt{cls()}, or \texttt{name()} live.
So it records the expression source code internally (\texttt{reads: \{key, code\}[]}) but does not pretend it can create closures.

The end-to-end test in \texttt{compiler/\_\_tests\_\_/compiler.e2e.test.ts} shows the intended integration:
the caller wires in the \texttt{read} closures manually.

\begin{lstlisting}[language=JavaScript]
const { block } = compileTSXToBlock(`(<div className={count()}>Count</div>)`)

const mounted = mountCompiledBlock(block, host, [
	{ key: 's0', read: () => count() },
])
\end{lstlisting}

The key behavioral contract in that test is worth repeating:

\begin{quote}
	extbf{Updating signals must never recreate the root element.}
\end{quote}

In other words, compilation is only useful if it reduces work to \emph{writes}, not \emph{mount/unmount}.

\section{Babel transform: how the output program is shaped}

The Babel plugin in \texttt{compiler/jsx/transform.ts} exists to solve the closure problem.
It runs inside your real source file, so it can emit \texttt{read: () => expr} in-place.

\subsection{The shape (as asserted by tests)}

The unit tests in \texttt{compiler/\_\_tests\_\_/babel-transform.test.ts} basically define the public contract.
The transform must emit:

\begin{itemize}
\item an import from \texttt{relax/hrbr} that includes \texttt{defineBlock} and \texttt{mountCompiledBlock}
\item a module-scope \texttt{const \_\_hrbr\_block\_N = defineBlock(...)}
\item component output that returns a \textbf{mount factory}: \texttt{host => \_\_hrbrBlock(host, def, slots)}
\item a wrapper function \texttt{\_\_hrbrBlock(host, def, slots)} that calls \texttt{mountCompiledBlock(def, host, slots)}
\end{itemize}

Here\'s the key part of the snapshot (reformatted as teaching code):

\begin{lstlisting}[language=JavaScript]
import { defineBlock, mountCompiledBlock } from "relax/hrbr";

const __hrbr_block_0 = defineBlock({
	templateHTML: "<section id=\"root\" class=\"\">\n...",
	slots: {
		"s0": { kind: "class", path: [] },
		"s1": { kind: "text", path: [3, 1] },
	}
});

const App = () => host => __hrbrBlock(host, __hrbr_block_0, [
	{ key: "s0", read: () => cls() },
	{ key: "s1", read: () => name() },
]);

function __hrbrBlock(host, def, slots) {
	return mountCompiledBlock(def, host, slots);
}
\end{lstlisting}

There are two deeply important design choices here:

\begin{enumerate}
\item \textbf{The \texttt{BlockDef} is module-scope data.}
	It should be created once, not per render.
\item \textbf{The mount factory is per render, and it closes over reads.}
	That\'s where \texttt{cls()} and \texttt{name()} are actually in scope.
\end{enumerate}

This pattern is very similar in spirit to Solid\'s compilation strategy (static DOM + fine-grained updates), though the runtime and APIs differ.
For a conceptual reference, see:
\href{https://www.solidjs.com/docs/latest\#compiling}{https://www.solidjs.com/docs/latest\#compiling}.

\subsection{Lane pragmas: compiler-to-scheduler hints without changing semantics}

Both compilers support a tiny \"lane pragma\" syntax:

\begin{lstlisting}[language=JavaScript]
<div
	className={/*@lane transition*/ cls()}
	data-x={/*@lane input*/ x()}
/>
\end{lstlisting}

The Babel test asserts that these land in slot metadata as \texttt{lane: "transition"} and \texttt{lane: "input"}.

Why it matters: when you start integrating compilation into an app/runtime scheduler, this gives you a way to nudge \emph{when} updates flush without changing \emph{what} updates mean.

\section{Debugging: when a slot path is wrong}

The most common failure mode when you build a compiler like this is:\
	extbf{the slot writes go to the wrong node.}

When that happens, you should immediately suspect one of these:

\begin{itemize}
\item you counted \texttt{children} instead of \texttt{childNodes}
\item your template string gained or lost whitespace (newlines/indentation)
\item you treated a void tag as having children (or forgot it closes itself)
\end{itemize}

\subsection*{Mini-exercise (tests-as-spec)}

Open \texttt{compiler/\_\_tests\_\_/babel-transform.test.ts} and find the inline snapshot where \texttt{path: [3, 1]}.
Then:

\begin{enumerate}
\item remove the whitespace/newlines in the template emitter (so the HTML becomes a single line),
\item update the snapshot, and
\item explain (in one sentence) why the path changed.
\end{enumerate}

This is the simplest way to internalize the \texttt{childNodes} rule.

The critical detail is that the compiler counts \emph{every child node} (elements \emph{and} text nodes), because \texttt{childNodes} includes text nodes. In \texttt{compiler/index.ts} this is represented by \texttt{PathState\{ childIndex \}}.

This is also why the compiler uses explicit sentinels for some dynamic expressions: it needs to ensure a real, countable node exists at the path.

\subsection*{Diagram: how slot paths are interpreted via \texttt{childNodes}}
\begin{center}
\begin{tikzpicture}[
	node distance=10mm and 14mm,
	box/.style={draw, rounded corners, align=left, inner sep=6pt, fill=RelaxLight},
	arrow/.style={-Latex, thick}
]
\node[box] (template) {\textbf{Compiler emits}\\\texttt{path = [3,1]}\\meaning: under template root,\\go to \texttt{childNodes[3]} then \texttt{childNodes[1]}};
\node[box, right=of template] (parse) {Runtime parses\\\texttt{templateHTML}\\into DOM tree};
\node[box, right=of parse] (walk) {Runtime resolves path\\\texttt{node = root}\\\texttt{for i in path: node = node.childNodes[i]}};
\node[box, below=of walk] (write) {Apply slot write\\(setText / setAttribute / etc.)};

\draw[arrow] (template) -- (parse);
\draw[arrow] (parse) -- (walk);
\draw[arrow] (walk) -- (write);
\end{tikzpicture}
\end{center}

This is why both compilers count \emph{text nodes} as children: the runtime uses \texttt{childNodes}, not \texttt{children}.

\subsection{Text expression handling: sentinel + slot path}

In \texttt{renderJSX()}, when a child is \texttt{JsxExpression}:

\begin{itemize}
\item the compiler marks \texttt{meta.hasDynamicStructure = true}
\item it allocates a slot key \texttt{sN}
\item it emits a unique sentinel string \texttt{\_\_hrbr\_sN\_\_} into the template
\item it records a \texttt{text} slot whose path points to that text node
\item it records the expression's source text as a ``read'' entry
\end{itemize}

This is a deliberately conservative design: any child expression could be a conditional or a list, so the prototype treats it as structure-changing.
The Babel compiler (later in this book) uses heuristics to allow some expressions as ``text-like'' without falling back.

\subsection{Void elements and template correctness}

The TS compiler treats HTML void tags specially via the \texttt{VOID\_TAGS} set.
For these tags it emits \texttt{<input ...>} rather than \texttt{<input ...></input>}.
This matters because the runtime parses \texttt{templateHTML} through \texttt{<template>} and relies on the browser's HTML parser to build a stable node tree. Producing non-HTML-canonical markup can shift the parsed tree and break slot paths.

\subsection{Attributes: literals vs expression slots}

The prototype supports:

\begin{itemize}
\item boolean attributes: \texttt{<div hidden />}
\item string literal attributes: \texttt{<div id="x" />}
\item expression attributes: \texttt{<div className=\{count()\} />}
\end{itemize}

Expression attributes allocate slot keys and record \texttt{kind} and \texttt{name}:

\begin{itemize}
\item \texttt{className} normalizes to \texttt{class}
\item \texttt{style} maps to a \texttt{style} slot
\item anything else becomes an \texttt{attr} slot
\end{itemize}

The compiler also records lane hints (per slot) and ``signal-like'' hints:

\begin{itemize}
\item \texttt{meta.lanesByKey}: a compiler-to-runtime hint used by schedulers in later phases
\item \texttt{meta.signalLikeByKey}: a heuristic that recognizes \texttt{foo()} style getters
\end{itemize}

Even when later phases choose to ignore (or reinterpret) these hints, keeping the meta separate from the structural template makes it possible to evolve scheduling policy without changing the fundamental block ABI.

\subsection{Fallback decision wrapper: \texttt{compileTSXOrFallback()}}

This wrapper is the ``policy layer'' on top of template extraction.
It answers one question: \emph{is it safe to run in block mode?}

At the top level, the TS prototype exposes a convenience API:
	exttt{compileTSXOrFallback(source)}.

Its semantics are intentionally boring:

\begin{lstlisting}[language=JavaScript]
const res = compileTSXToBlock(source)

if (res.meta.hasDynamicStructure) {
	return { kind: 'fallback', reason: 'dynamic-structure', meta: res.meta }
}

return { kind: 'block', ...res }
\end{lstlisting}

This is important because it makes ``fallback'' a first-class output, not an exception.
The compiler is allowed to say ``I don't know''.

\section{Dynamic structure detection (what triggers fallback)}

Now we'll make the fuzzy idea from earlier (\emph{structure vs values}) precise.
The easiest way is to turn the implementation into a checklist.

\subsection{Decision table: TS prototype vs Babel transform}

Relax currently has a \emph{conservative} detector in the TS prototype and a \emph{slightly more permissive} detector in the Babel transform.
They differ because Babel has access to more information (and can safely reject with fallback by returning \texttt{null} from template building).

\begin{center}
\begin{tabular}{|p{0.34\textwidth}|p{0.29\textwidth}|p{0.29\textwidth}|}
\hline
	extbf{Input pattern} & \textbf{TS prototype (\texttt{compiler/index.ts})} & \textbf{Babel transform (\texttt{compiler/jsx/transform.ts})}\\
\hline
Child expression: \texttt{<p>Hello \{expr\}</p>} & Marks \texttt{meta.hasDynamicStructure = true} (always) & Allows only \emph{text-like} expressions; otherwise returns \texttt{null} to route to fallback \\
\hline
Conditional/logical children: \texttt{cond ? <A/> : <B/>}, \texttt{cond \&\& <A/>} & Treated as dynamic structure (child expression) & Treated as ``not text-like'' (\texttt{ConditionalExpression}/\texttt{LogicalExpression}); routes to fallback \\
\hline
Non-intrinsic tags: \texttt{<Component/>}, \texttt{<Foo.Bar/>} & Marks dynamic structure when tag is not lowercase & Template builder returns \texttt{null} when tag is not a lowercase intrinsic \\
\hline
Event handlers: \texttt{onClick=\{fn\}} & Marks dynamic structure (events not representable in static HTML) & Immediately returns \texttt{null} (block mode not supported for events yet) \\
\hline
Spread attributes: \texttt{<div \{...props\} />} & Marks dynamic structure (and throws in strict template mode) & Returns \texttt{null} (not supported in block template) \\
\hline
Void tags: \texttt{<input/>}, \texttt{<img/>} & Emitted with correct HTML void syntax via \texttt{VOID\_TAGS} & Emitted with a regex-based void-tag check \\
\hline
\end{tabular}
\end{center}

This table is the practical meaning of \emph{``dynamic structure''} in Relax today.

\subsection{Babel's \texttt{isTextLikeExpr}: the tiny heuristic that buys you a lot}

The Babel transform has a helper called \texttt{isTextLikeExpr(expr)}.
The idea is simple: if a child expression could evaluate to an element/fragment/array, then the template structure isn't stable.
But if it's always a value that can be written into \texttt{textContent}, we can keep block mode.

From the scaffold (summarized):

\begin{lstlisting}[language=JavaScript]
function isTextLikeExpr(expr) {
	if (expr is JSXElement or JSXFragment) return false
	if (expr is ConditionalExpression) return false
	if (expr is LogicalExpression) return false
	return true
}
\end{lstlisting}

Notice what this buys us:

\begin{itemize}
\item \textbf{Text updates stay cheap.} Most UIs do a lot of ``render some string'' work.
\item \textbf{You can still escape hatch.} The moment an expression might be structural, it falls back.
\end{itemize}

This is a deliberately small heuristic.
In later chapters, you'll see how it grows, and how tests keep it honest.

\section{Babel compilation mechanics: how one JSX element becomes a block}

So far we've looked at the \emph{shape} of the output.
Now we'll build a mental model for the \emph{mechanics} of how the Babel transform creates that output.

The transform in \texttt{compiler/jsx/transform.ts} does three jobs:

\begin{enumerate}
\item build a static HTML template string
\item allocate slots + paths for the dynamic parts
\item emit a mount factory that closes over live reads
\end{enumerate}

If you understand those three, the whole transform stops feeling magical.

\subsection{Step 1: pick a stable root (or give up)}

The transform only wants to compile stable, intrinsic HTML tags like \texttt{div}, \texttt{section}, and \texttt{button}.
If the tag is not a lowercase intrinsic (for example \texttt{<Component/>} or \texttt{<Foo.Bar/>}), there is no obvious static HTML representation.
So the template builder returns \texttt{null}, and the plugin routes to fallback.

Why is this the right default?

\begin{itemize}
\item In JSX, a capitalized tag means ``call some function and let it return anything''.
\item That means the compiler can't promise stable structure.
\end{itemize}

Block mode only works when the compiler can promise that \emph{the same DOM nodes exist forever}.

\subsection{Step 2: render a string that the browser will parse deterministically}

Think of compilation as emitting a \texttt{<template>} string.
Your compiler isn't building DOM; the browser will.
So your job is just to produce a string that will be parsed into a stable tree.

That implies two constraints:

\begin{itemize}
\item \textbf{Void tags matter.} \texttt{<input>} can't accidentally become \texttt{<input></input>}.
\item \textbf{Whitespace matters.} Newlines/indentation become text nodes and shift \texttt{childNodes} indices.
\end{itemize}

This is why the snapshot tests in \texttt{compiler/\_\_tests\_\_/babel-transform.test.ts} are so valuable.
They lock down not just the slot paths, but the \emph{exact} template string (including newlines).

\subsection{Step 3: allocate slots: ``static HTML with holes''}

When the transform sees something dynamic, it does \emph{not} generate DOM operations directly.
It just records metadata that the runtime knows how to apply later.

You can think of slot allocation as making a table that answers:

\begin{quote}
When the runtime mounts this template, \emph{where} is the thing to update, and \emph{how} should it be updated?
\end{quote}

Examples you should expect as you read the code:

\begin{itemize}
\item \textbf{Attribute slot:} dynamic attribute value (\texttt{kind: 'attr'}).
\item \textbf{Class slot:} \texttt{className=\{...\}} normalized to \texttt{class} (\texttt{kind: 'class'}).
\item \textbf{Text slot:} dynamic text content inside an element (\texttt{kind: 'text'}).
\end{itemize}

Every slot includes a \texttt{path}, which is just a frozen version of the \texttt{childNodes} walk we discussed earlier.

\subsection{Step 4: emit reads as closures (the key Babel advantage)}

The TS prototype can't produce \texttt{read: () => expr} because it doesn't know which variables exist.
But Babel is transforming your actual file, so it can emit closures that reference \texttt{cls()}, \texttt{name()}, or any local binding.

The key idea is:

\begin{quote}
	extbf{All dynamic values are read through closures passed to \texttt{mountCompiledBlock()}.}
\end{quote}

This means compilation doesn't change your state model.
Signals are still signals.
The compiler just gives the runtime a faster way to apply their values.

\subsection{Step 5: fallback routing is just ``emit a different mount factory''}

In the Babel plugin, ``fallback'' isn't a special runtime mode.
It's just a different function call in the generated output:

\begin{itemize}
\item block mode: emit \texttt{mountCompiledBlock(def, host, compiledSlots)}
\item fallback mode: emit \texttt{mountFallback(host, () => <JSX...>)} (shape varies as the scaffold evolves)
\end{itemize}

The important part is the architecture: both modes still return a mount function.
So the caller doesn't care which one it got.

For conceptual background on why many UI frameworks separate ``render'' from ``commit'' (build a description then apply it), see:
\href{https://react.dev/learn/render-and-commit}{https://react.dev/learn/render-and-commit}.

\section{Validation contract: how to change the compiler safely}

Compilers are notoriously easy to break with tiny edits.
Relax deals with this by treating tests as the spec.
If you're modifying the compiler, keep these two mental buckets:

\begin{itemize}
\item \textbf{Behavioral contracts:} integration tests that assert DOM stability.
\item \textbf{Shape contracts:} snapshot tests that assert exact emitted output.
\end{itemize}

\subsection{Spec-to-test mapping (start here when debugging)}

\begin{itemize}
\item \textbf{"Don't recreate root DOM" contract:} \texttt{compiler/\_\_tests\_\_/compiler.e2e.test.ts}
\item \textbf{"Exact codegen shape" contract (imports, defineBlock, wrapper):} \texttt{compiler/\_\_tests\_\_/babel-transform.test.ts}
\item \textbf{"Path semantics" contract (whitespace affects \texttt{childNodes}):} also \texttt{babel-transform.test.ts} snapshots
\end{itemize}

When you're unsure whether a change is safe, prefer adding a tiny unit test over adding another heuristic.
Heuristics get clever; tests keep you honest.

Now we\'ll make the fuzzy idea from earlier (\emph{structure vs values}) precise.
The easiest way is to turn the implementation into a checklist.

\subsection{Decision table: TS prototype vs Babel transform}

Relax currently has a \emph{conservative} detector in the TS prototype and a \emph{slightly more permissive} detector in the Babel transform.
They differ because Babel has access to more information (and can safely reject with fallback by returning \texttt{null} from template building).

\begin{center}
\begin{tabular}{|p{0.34\textwidth}|p{0.29\textwidth}|p{0.29\textwidth}|}
\hline
	extbf{Input pattern} & \textbf{TS prototype (\texttt{compiler/index.ts})} & \textbf{Babel transform (\texttt{compiler/jsx/transform.ts})}\\
\hline
Child expression: \texttt{<p>Hello \{expr\}</p>} & Marks \texttt{meta.hasDynamicStructure = true} (always) & Allows only \emph{text-like} expressions; otherwise returns \texttt{null} to route to fallback \\
\hline
Conditional/logical children: \texttt{cond ? <A/> : <B/>}, \texttt{cond \&\& <A/>} & Treated as dynamic structure (child expression) & Treated as \\"not text-like\\" (\texttt{ConditionalExpression}/\texttt{LogicalExpression}); routes to fallback \\
\hline
Non-intrinsic tags: \texttt{<Component/>}, \texttt{<Foo.Bar/>} & Marks dynamic structure when tag is not lowercase & Template builder returns \texttt{null} when tag is not a lowercase intrinsic \\
\hline
Event handlers: \texttt{onClick=\{fn\}} & Marks dynamic structure (events not representable in static HTML) & Immediately returns \texttt{null} (block mode not supported for events yet) \\
\hline
Spread attributes: \texttt{<div \{...props\} />} & Marks dynamic structure (and throws in strict template mode) & Returns \texttt{null} (not supported in block template) \\
\hline
Void tags: \texttt{<input/>}, \texttt{<img/>} & Emitted with correct HTML void syntax via \texttt{VOID\_TAGS} & Emitted with a regex-based void-tag check \\
\hline
\end{tabular}
\end{center}

This table is the practical meaning of \emph{\"dynamic structure\"} in Relax today.

\subsection{Babel\'s \texttt{isTextLikeExpr}: the tiny heuristic that buys you a lot}

The Babel transform has a helper called \texttt{isTextLikeExpr(expr)}.
The idea is simple: if a child expression could evaluate to an element/fragment/array, then the template structure isn\'t stable.
But if it\'s always a value that can be written into \texttt{textContent}, we can keep block mode.

From the scaffold (summarized):

\begin{lstlisting}[language=JavaScript]
function isTextLikeExpr(expr) {
	if (expr is JSXElement or JSXFragment) return false
	if (expr is ConditionalExpression) return false
	if (expr is LogicalExpression) return false
	return true
}
\end{lstlisting}

Notice what this buys us:

\begin{itemize}
\item \texttt{<span>{count()}</span>} can be compiled as a text slot.
\item \texttt{<span>{count() > 0 ? <b/> : null}</span>} must go to fallback.
\end{itemize}

Even though this is conservative, it gets you most of the performance wins for the common case (text/attr/class/style bindings).

\section{How fallback routing looks in the Babel output (scaffold)}

When \texttt{buildTemplateFromJsxElement()} returns \texttt{null}, the plugin constructs a minimal \texttt{mountFallback()} call.
This is still early-stage code, but it\'s worth understanding the design:

\begin{itemize}
\item block mode emits a \texttt{BlockDef} and a tight \texttt{CompiledSlot[]} list
\item fallback emits a \emph{render function} that returns reconciliation specs (create/patch)
\end{itemize}

That means fallback is not \"VDOM\" in the React sense.
It\'s a small structural reconciler that still aims to avoid recreating nodes whenever possible.

If you want a conceptual baseline for why compilers do this \"static template + runtime writes\" split, React\'s docs on render/commit are a good mental model:
\href{https://react.dev/learn/render-and-commit}{https://react.dev/learn/render-and-commit}.

\subsection*{Spec-to-test mapping (what to read next)}

When you\'re debugging compiler behavior, treat these tests as your spec:

\begin{itemize}
\item \textbf{No DOM recreation on updates}: \texttt{compiler/\_\_tests\_\_/compiler.e2e.test.ts}
\item \textbf{Output program shape}: \texttt{compiler/\_\_tests\_\_/babel-transform.test.ts}
\end{itemize}

If a behavior isn't asserted there yet, it's not stable.

\texttt{compileTSXOrFallback()} runs \texttt{compileTSXToBlock()} and then returns:

\begin{itemize}
\item \texttt{\{ kind: 'fallback', reason: 'dynamic-structure' \}} if \texttt{meta.hasDynamicStructure} is true
\item otherwise \texttt{\{ kind: 'block', ...result \}}
\end{itemize}

This is the high-level policy the real Babel transform is evolving toward.

Note the important nuance: \texttt{compileTSXToBlock()} still returns a \texttt{BlockDef} even when \texttt{meta.hasDynamicStructure} is true. That mirrors how the Babel plugin can still hoist block defs for other expressions in the same file even if one subtree needs fallback.

\section{Babel transform scaffold (real-world integration)}

The Babel plugin entrypoint is exported as \texttt{hrbrJsxTransform} from \texttt{compiler/jsx/index.ts} and re-exported from \texttt{compiler/index.ts}.
The implementation lives in \texttt{compiler/jsx/transform.ts}.

\subsection{Roles of the Babel plugin}

Compared to the TS prototype, the Babel plugin has additional responsibilities:

\begin{itemize}
\item Operate on full program files, not just an isolated expression string.
\item Inject runtime imports: \texttt{defineBlock}, \texttt{mountCompiledBlock}, and optionally \texttt{mountFallback}.
\item Generate wrapper helpers (e.g. \texttt{\_\_hrbrBlock}) that connect emitted data literals to runtime calls.
\item Derive stable slot keys in dev mode (optionally) so diffs are readable.
\end{itemize}

\subsection{``Block mode'' is best-effort: return \texttt{null} to fall back}

The internal function \texttt{buildTemplateFromJsxElement()} tries to convert a JSX element into:

\begin{itemize}
\item \texttt{templateHTML}
\item \texttt{slots} map
\item \texttt{compiledSlots} array literal (with \texttt{read: () => expr})
\end{itemize}

If it sees unsupported patterns (events, dynamic elements, structural expressions), it returns \texttt{null} and marks that fallback imports are needed.

\subsection{Babel transform: what ``best-effort block mode'' means in code}

In \texttt{compiler/jsx/transform.ts}, block mode is essentially:

\begin{itemize}
\item Try to turn a \texttt{JSXElement} into \texttt{\{ templateHTML, slots, compiledSlots \}} using \texttt{buildTemplateFromJsxElement()}.
\item If any unsupported structure appears, return \texttt{null}.
\item The caller then chooses a fallback mount factory instead of a compiled-block mount factory.
\end{itemize}

The current transform is strict on purpose (e.g. it rejects events and spreads in block mode), because the project is trying to keep the block subset crisp while the fallback runtime is still evolving.

\subsection{Slot key strategies: sequential vs stable (dev mode)}

The Babel compiler contains two different slot key allocation styles:

\begin{itemize}
\item \textbf{Default (prod-like):} sequential keys \texttt{s0, s1, ...}.
\item \textbf{Dev mode:} location-based keys like \texttt{s\_text\_12:18} derived from source positions.
\end{itemize}

Location-based keys are helpful because they reduce ``slot churn'' in diffs and make debugging invalid-path errors more readable.

\subsection{Lane pragmas (compiler-to-scheduler hint)}

The Babel plugin recognizes comment pragmas in expression containers:

\begin{quote}
\texttt{className=\{/*@lane transition*/ cls()\}}
\end{quote}

It extracts the lane and stores it on the slot metadata.
Later chapters will connect this to the deterministic scheduler.

Concretely, \texttt{extractLanePragmaFromExpressionContainer()} scans leading comments for a pattern like:\\
	texttt{@lane transition}.
When found, the compiler attaches \texttt{lane: "transition"} to the slot metadata object.

\section{Tests-as-spec}

The compiler has three layers of tests:

\begin{itemize}
\item Unit compilation tests (shape of emitted template/slots): \texttt{compiler/\_\_tests\_\_/compiler.test.ts}
\item Snapshot tests for stability: \texttt{compiler/\_\_tests\_\_/compiler.snapshot.test.ts}
\item End-to-end tests that mount the result in jsdom:
  \begin{itemize}
  \item \texttt{compiler/\_\_tests\_\_/compiler.e2e.test.ts} (TS prototype -> mountCompiledBlock)
  \item \texttt{compiler/\_\_tests\_\_/babel-transform.e2e-jsdom.test.ts} (Babel output -> evaluate -> mountCompiledBlock)
  \end{itemize}
\end{itemize}

In the rest of this section, we'll treat the tests as a specification of the compiler/runtime boundary.
We won't run them here, but reading them tells you what changes are allowed without breaking the public contract.

\subsection{Unit tests (TS): compile + mount sanity}


	extbf{File:} \texttt{compiler/\_\_tests\_\_/compiler.test.ts}

This test anchors the smallest happy path:

\begin{itemize}
\item \texttt{compileTSXToBlock()} accepts a simple intrinsic element with a dynamic attribute expression.
\item The caller supplies \texttt{CompiledSlot\{ key: 's0', read: () => name() \}}.
\item \texttt{mountCompiledBlock()} produces DOM matching the template, and the dynamic slot lands in the right place.
\end{itemize}

One subtle detail: the test disposes the mounted block via \texttt{mounted.dispose()}. This keeps the contract clear: compiled mounts are still reactive, and disposal is part of the lifecycle.

\subsection{TS end-to-end: compile -> mount -> update without recreating root}

In \texttt{compiler/\_\_tests\_\_/compiler.e2e.test.ts}:

\begin{enumerate}
\item Create a signal: \texttt{const [count, setCount] = createSignal(1)}.
\item Compile TSX: \texttt{(<div className=\{count()\}>Count</div>)}.
\item Mount via \texttt{mountCompiledBlock(block, host, \{ key: 's0', read: () => count() \})}.
\item Update the signal, wait for the scheduler (microtask + \texttt{setTimeout(0)}).
\item Assert the root element identity is preserved and only the class changed.
\end{enumerate}

This is the key end-to-end contract for compiled blocks.

This test also documents a scheduling assumption: it comments that HRBR scheduling uses \texttt{setTimeout(0)} by default in tests, so updates require waiting at least one timer tick.

\subsection{Babel transform unit tests: shape of emitted program}


	extbf{File:} \texttt{compiler/\_\_tests\_\_/babel-transform.test.ts}

These tests assert that the plugin is producing the intended \emph{program shape}:

\begin{itemize}
\item it injects an import from the configured runtime path (default \texttt{relax/hrbr})
\item it references \texttt{defineBlock} and \texttt{mountCompiledBlock}
\item it hoists a \texttt{const \_\_hrbr\_block\_0 = defineBlock(...)}
\item it emits a wrapper mount function of the form \texttt{host => \_\_hrbrBlock(host, def, slots)}
\end{itemize}

The snapshot test in this file is especially important: it pins down the exact emitted structure (templateHTML string, slots map, and compiledSlots array with \texttt{read: () => expr}).

The suite also asserts two ``developer experience'' guarantees:

\begin{itemize}
\item lane pragmas are propagated to slot metadata
\item dev mode can emit location-based keys
\end{itemize}

\subsection{Babel transform errors: the intentional block-mode subset}


	extbf{File:} \texttt{compiler/\_\_tests\_\_/babel-transform.errors.test.ts}

This suite documents what is currently \emph{rejected} in block mode:

\begin{itemize}
\item event handlers (e.g. \texttt{onClick})
\item spread attributes (\texttt{<div \{...props\} />})
\item fragments (\texttt{<>...</>})
\item non-intrinsic tags (components like \texttt{<Foo />})
\end{itemize}

Even if future versions route these to fallback instead of throwing, the existence of these tests is useful: it makes the current compilation subset explicit and prevents accidental expansion that would silently generate incorrect templates.

\subsection{Babel transform E2E (jsdom): evaluate emitted literals + mount}


	extbf{File:} \texttt{compiler/\_\_tests\_\_/babel-transform.e2e-jsdom.test.ts}

This test checks something subtle but crucial: the compiler output is not just syntactically valid---its emitted \emph{data} can actually be evaluated and mounted.

Notable aspects:

\begin{itemize}
\item It extracts the object literal passed to \texttt{defineBlock(...)} from the emitted text.
\item It extracts the compiledSlots array literal from the wrapper call.
\item It evaluates those two pieces with \texttt{new Function(...)} while providing \texttt{cls()} and \texttt{name()} from the test scope.
\item It mounts using \texttt{mountCompiledBlock} and asserts identity preservation across updates.
\end{itemize}

This approximates how ``real code'' will behave: the compile-time output must work with runtime scope (signals) and produce stable DOM.

\section{Worked example: derive \texttt{templateHTML} and slot paths from the snapshot}

Now we\'ll do the most valuable \emph{derive-it-by-hand} exercise: take the golden snapshot in
	exttt{compiler/\_\_tests\_\_/babel-transform.test.ts} and derive the template and slot paths by hand.
If you can do this once, debugging a compiler becomes much less mysterious.

\subsection{The input JSX (what the author writes)}

The snapshot test compiles this input:

\begin{lstlisting}[language=JavaScript]
const App = () => (
	<section id="root" className={cls()}>
		<h1>Title</h1>
		<p>Hello {name()}</p>
	</section>
)
\end{lstlisting}

At a high level, the compiler must split this into:

\begin{itemize}
\item \textbf{Structure:} a static HTML template string.
\item \textbf{Values:} a list of dynamic slots with deterministic DOM paths.
\end{itemize}

\subsection{Step 1: the static template string (structure + placeholders)}

In the normalized snapshot, the transform emits a \texttt{defineBlock({ templateHTML, slots })} where \texttt{templateHTML} is a single string.
It contains the \texttt{<section>} with \texttt{id="root"}, plus an empty \texttt{class=""} placeholder.

In the snapshot, the template also includes indentation and newlines.
That\'s not just cosmetic: when the runtime parses HTML, whitespace becomes text nodes.

Why the placeholders matter:

\begin{itemize}
\item \texttt{className=\{cls()\}} becomes \texttt{class=""} in the template so the runtime can later write the real attribute value.
\item \texttt{Hello \{name()\}} becomes \texttt{Hello  } in the template (leaving a stable ``spot'' for the text slot).
\end{itemize}

If you ever see a compiler try to ``just concatenate strings'' at runtime, this is the key difference: HRBR wants the template to be parsed once, and then updated via direct DOM writes.

\subsection{Step 2: the slot declarations (what to write + where to write it)}

The same snapshot asserts two slot declarations:

\begin{itemize}
\item \texttt{s0}: \texttt{kind: "class"}, \texttt{path: []}
\item \texttt{s1}: \texttt{kind: "text"}, \texttt{path: [3, 1]}
\end{itemize}

Think of \texttt{path} as a precomputed set of array indexes into \texttt{childNodes}.
No search. No diff. Just indexing.

\subsection{Step 3: compute \texttt{path: [3, 1]} by hand (the whitespace trap)}

We\'ll compute \texttt{[3, 1]} mechanically.

First hop: \texttt{3} means we\'re selecting \texttt{section.childNodes[3]}.
Because the template contains whitespace between elements, the \texttt{childNodes} list includes \emph{text nodes}:

\begin{enumerate}
\item \texttt{childNodes[0]}: text (newline/indent before \texttt{<h1>})
\item \texttt{childNodes[1]}: \texttt{<h1>}
\item \texttt{childNodes[2]}: text (newline/indent between \texttt{</h1>} and \texttt{<p>})
\item \texttt{childNodes[3]}: \texttt{<p>}
\item \texttt{childNodes[4]}: text (newline/indent before \texttt{</section>})
\end{enumerate}

So hop 1 lands on the \texttt{<p>}.

Second hop: \texttt{1} means \texttt{p.childNodes[1]}.
The \texttt{<p>} contains static text (\texttt{Hello }) plus a dynamic text slot.
The compiler\'s job is to ensure there\'s a stable text node representing the dynamic value.
So we get:

\begin{itemize}
\item \texttt{p.childNodes[0]}: text node for \texttt{"Hello "}
\item \texttt{p.childNodes[1]}: text node that will receive \texttt{name()} updates
\end{itemize}

That\'s what \texttt{path: [3, 1]} means: ``the second text node inside the paragraph.''

This is the most important compiler invariant in HRBR: if you change the template emitter and accidentally add/remove whitespace, you can shift every downstream path.

\subsection{Step 4: compiled slot readers (how values cross the boundary)}

The final piece the snapshot asserts is the compiled slot array literal:

\begin{lstlisting}[language=JavaScript]
const App = () => host => __hrbrBlock(host, __hrbr_block_0, [
	{ key: "s0", read: () => cls() },
	{ key: "s1", read: () => name() },
])
\end{lstlisting}

This is where Babel shines: \texttt{read} is a closure in the output program.
There\'s no serialization format for expressions.
The emitted code just calls the original functions.

\subsection{Debugging checklist (tests-first)}

When a compiled block updates the wrong node, treat it like a deterministic math error.
Use the tests to narrow it down:

\begin{enumerate}
\item Check the snapshot: \texttt{compiler/\_\_tests\_\_/babel-transform.test.ts}.
  If your change modified whitespace in \texttt{templateHTML}, you probably shifted paths.
\item Recompute the \texttt{childNodes} list by hand and see if the asserted \texttt{path} still lands on the intended node.
\item Run the E2E that actually mounts the evaluated literals:
  	exttt{compiler/\_\_tests\_\_/babel-transform.e2e-jsdom.test.ts}.
\end{enumerate}

If you get into the habit of adding a snapshot test first, you\'ll catch most path regressions before they reach runtime.

\subsection{Babel fallback E2E (jsdom): structural expression routes to \texttt{mountFallback}}


	extbf{File:} \texttt{compiler/\_\_tests\_\_/babel-transform.fallback.e2e-jsdom.test.ts}

This suite pins down the fallback routing policy:

\begin{itemize}
\item A conditional expression in children is treated as \textbf{dynamic structure} (unsafe for block mode).
\item The transform output must include \texttt{mountFallback}.
\item The emitted module is evaluated, and calling \texttt{App()} returns a mount function.
\item Updates are reactive: signal changes trigger DOM updates without throwing.
\end{itemize}

The test includes a comment that the MVP fallback spec doesn't yet clear conditional text in all cases; the contract being asserted is deliberately smaller: the pipeline compiles, mounts, and updates safely.

\section{Where we go next}

Chapter~\ref{ch:babel-plugin} will zoom in on the Babel visitor model and show how the plugin emits the full final shape that user code will run.
We will also connect the ``fail block mode'' paths to the fallback runtime and reconciler described earlier.

\section{Online resources}

\begin{itemize}
\item TypeScript compiler API: parsing source and working with AST nodes:\\
\href{https://github.com/microsoft/TypeScript/wiki/Using-the-Compiler-API}{TypeScript Wiki: Using the Compiler API}.

\item Babel plugin authoring and visitor model:\\
\href{https://babeljs.io/docs/en/plugins}{Babel docs: plugins},\\
\href{https://babeljs.io/docs/en/babel-plugin-api}{Babel docs: plugin API}.

\item JSX specification details (useful when deciding what constitutes ``dynamic structure''):\\
\href{https://facebook.github.io/jsx/}{JSX spec (historical reference)}.

\item Why HTML parsing matters for template correctness (void elements, tree building, etc.):\\
\href{https://developer.mozilla.org/en-US/docs/Web/HTML/Element}{MDN: HTML elements reference}.

\item Source maps basics (relevant to dev mode and debugging transforms):\\
\href{https://sourcemaps.info/spec.html}{Source Map Revision 3 Proposal}.
\end{itemize}
